\documentclass[11pt,a4paper]{article}
\usepackage[T1]{fontenc}
\usepackage[utf8]{inputenc}
\usepackage{amsmath,amssymb,amsthm}
\usepackage[margin=1in]{geometry}
\usepackage[colorlinks=true,linkcolor=blue,urlcolor=blue]{hyperref}
\theoremstyle{plain}
\newtheorem{theorem}{Theorem}
\newtheorem{lemma}[theorem]{Lemma}
\newtheorem{proposition}[theorem]{Proposition}
\newtheorem{corollary}[theorem]{Corollary}
\theoremstyle{definition}
\newtheorem{remark}[theorem]{Remark}

\title{The empirical recurrence for OEIS A209793, proved by transfer matrix}
\author{Adrian Perez Fontelles\\ \small Independent researcher}
\date{31 August 2026}
\begin{document}
\maketitle

\begin{abstract}
OEIS A209793 counts $(n+1)\times6$ arrays over $\{0,\dots,2\}$ subject to a condition imposed on
every $2\times2$ subblock, and carries an empirical recurrence of order $70$ contributed
by R.~H.~Hardin. It is true, and it is not an empirical matter at all. The condition is local
to each subblock, so the rows of an admissible array are the vertices of a finite digraph
on $S=729$ vertices whose edges are the admissible consecutive pairs, and the count is a
number of walks: $a(n)=\tfrac1{2}\mathbf{1}^{\!\top}M^{n}\mathbf{1}$. Writing $q$
for the characteristic polynomial of the conjectured recurrence, the recurrence holds for all
$n$ exactly when $\mathbf{1}^{\!\top}M^{j}q(M)\mathbf{1}=0$ for every $j\ge0$, and the
Cayley--Hamilton theorem bounds that to $j<S$. It is therefore a finite computation in exact
integer arithmetic, and it comes out zero.
\end{abstract}

\noindent\small 2020 Mathematics Subject Classification. 05A15, 05B45, 15B36.\normalsize

\section{The sequence and the conjecture}

OEIS A209793 is ``Half the number of (n+1)X6 0..2 arrays with every 2X2 subblock having exactly one duplicate clockwise edge difference.''. It has offset $1$ and begins
\[
63300, 11761044, 2167312950, 400252237467, 73890766377408, 13642480312938969, 2518750352620857312, 465030007730370004089,\ \dots
\]
The entry carries the following comment:
\begin{quote}\small
Empirical: a(n) = 117*a(n-1) +16046*a(n-2) -461310*a(n-3) -41690248*a(n-4) +838185919*a(n-5) +37263733439*a(n-6) -765694604535*a(n-7) -13382141786885*a(n-8) +321364657017708*a(n-9) +2098104217159967*a(n-10) -69241805735327928*a(n-11) -128517018078543576*a(n-12) +8709129115010514837*a(n-13) -3468494802113380798*a(n-14) -700368603245220778907*a(n-15) +1097918366326791986076*a(n-16) +38249881277019367763149*a(n-17) -82885453957584747249841*a(n-18) -1477679050506703199996840*a(n-19) +3582295289065080645777908*a(n-20) +41539135165726319546746165*a(n-21) -101993734883143731283690081*a(n-22) -866781561552403343006767069*a(n-23) +2018826395772209491813887071*a(n-24) +13611750304478820781529014134*a(n-25) -28520372689086068572041492757*a(n-26) -162280092698848694950491324828*a(n-27) +291520634476059290564320104306*a(n-28) +1475295967406701632971876156473*a(n-29) -2169961121048984327488204378690*a(n-30) -10236687623863144159526172235403*a(n-31) +11783376804855161753282014450164*a(n-32) +54141369111431768097138884969134*a(n-33) -46589591407728596256741674132830*a(n-34) -217783303537298527601126461642677*a(n-35) +133306781440045365865140258465122*a(n-36) +664827799503260889733752180032703*a(n-37) -272365583756803560888165694739905*a(n-38) -1537257174407032961821202515767658*a(n-39) +385978218317055429280685779706059*a(n-40) +2686679225458054392142060493484229*a(n-41) -352927939490979467614872624551983*a(n-42) -3538593730688886357595470024571200*a(n-43) +158693762106843672238602958751326*a(n-44) +3496707953192547440599945751571573*a(n-45) +48384480903142552673396182803350*a(n-46) -2575506021288291072726617784495748*a(n-47) -127454659502298940357656903692957*a(n-48) +1400915090693069502783915455752061*a(n-49) +93099985963220323728526563241857*a(n-50) -555552327597734924323466692919042*a(n-51) -38316357184343791656582840644779*a(n-52) +157817720300615851244077379672100*a(n-53) +9697916680878993169585799823309*a(n-54) -31356941661649796672271184584347*a(n-55) -1525766261316131112529829339136*a(n-56) +4221660875174791218751179503189*a(n-57) +149433372546777023522642188062*a(n-58) -369752601975483264557434615861*a(n-59) -9422642718209654494532713305*a(n-60) +20000594578056871600802948792*a(n-61) +401829888361187527123789983*a(n-62) -619317950682081103026756420*a(n-63) -11022616921719495980402767*a(n-64) +9643967015220886173644550*a(n-65) +146861428770588720976758*a(n-66) -57353884938704235337372*a(n-67) -554801588254055287096*a(n-68) +40287240447192362128*a(n-69) +468254852850305856*a(n-70)
\end{quote}
As of the ``Last modified'' line on the live entry (Oct 03 2025 15:29:50, revision 9) the statement is
still recorded as empirical, and nothing on the entry records it as settled either way.

Written out, the claim is
\begin{equation}\label{eq:conj}
a(n)\;=\;117\,a(n-1) + 16046\,a(n-2) - 461310\,a(n-3) - 41690248\,a(n-4) + 838185919\,a(n-5) + 37263733439\,a(n-6) - 765694604535\,a(n-7) - 13382141786885\,a(n-8) + 321364657017708\,a(n-9) + 2098104217159967\,a(n-10) - 69241805735327928\,a(n-11) - 128517018078543576\,a(n-12) + 8709129115010514837\,a(n-13) - 3468494802113380798\,a(n-14) - 700368603245220778907\,a(n-15) + 1097918366326791986076\,a(n-16) + 38249881277019367763149\,a(n-17) - 82885453957584747249841\,a(n-18) - 1477679050506703199996840\,a(n-19) + 3582295289065080645777908\,a(n-20) + 41539135165726319546746165\,a(n-21) - 101993734883143731283690081\,a(n-22) - 866781561552403343006767069\,a(n-23) + 2018826395772209491813887071\,a(n-24) + 13611750304478820781529014134\,a(n-25) - 28520372689086068572041492757\,a(n-26) - 162280092698848694950491324828\,a(n-27) + 291520634476059290564320104306\,a(n-28) + 1475295967406701632971876156473\,a(n-29) - 2169961121048984327488204378690\,a(n-30) - 10236687623863144159526172235403\,a(n-31) + 11783376804855161753282014450164\,a(n-32) + 54141369111431768097138884969134\,a(n-33) - 46589591407728596256741674132830\,a(n-34) - 217783303537298527601126461642677\,a(n-35) + 133306781440045365865140258465122\,a(n-36) + 664827799503260889733752180032703\,a(n-37) - 272365583756803560888165694739905\,a(n-38) - 1537257174407032961821202515767658\,a(n-39) + 385978218317055429280685779706059\,a(n-40) + 2686679225458054392142060493484229\,a(n-41) - 352927939490979467614872624551983\,a(n-42) - 3538593730688886357595470024571200\,a(n-43) + 158693762106843672238602958751326\,a(n-44) + 3496707953192547440599945751571573\,a(n-45) + 48384480903142552673396182803350\,a(n-46) - 2575506021288291072726617784495748\,a(n-47) - 127454659502298940357656903692957\,a(n-48) + 1400915090693069502783915455752061\,a(n-49) + 93099985963220323728526563241857\,a(n-50) - 555552327597734924323466692919042\,a(n-51) - 38316357184343791656582840644779\,a(n-52) + 157817720300615851244077379672100\,a(n-53) + 9697916680878993169585799823309\,a(n-54) - 31356941661649796672271184584347\,a(n-55) - 1525766261316131112529829339136\,a(n-56) + 4221660875174791218751179503189\,a(n-57) + 149433372546777023522642188062\,a(n-58) - 369752601975483264557434615861\,a(n-59) - 9422642718209654494532713305\,a(n-60) + 20000594578056871600802948792\,a(n-61) + 401829888361187527123789983\,a(n-62) - 619317950682081103026756420\,a(n-63) - 11022616921719495980402767\,a(n-64) + 9643967015220886173644550\,a(n-65) + 146861428770588720976758\,a(n-66) - 57353884938704235337372\,a(n-67) - 554801588254055287096\,a(n-68) + 40287240447192362128\,a(n-69) + 468254852850305856\,a(n-70)
\end{equation}
for all $n>69$.

\section{The condition is local, so the rows form a digraph}

Write an array of the shape in question as its list of rows
$r^{(1)},r^{(2)},\dots$, each an element of
$V=\{0,1,\dots,2\}^{6}$, so $|V|=S=729$. A $2\times2$ subblock of the array
occupies two consecutive rows and two consecutive columns; if the two rows are $r$
and $s$ (in that order) and the two columns are numbered $j,j+1$,
the subblock is
\[
\begin{pmatrix}\alpha&\beta\\\gamma&\delta\end{pmatrix}=\begin{pmatrix}r_j&r_{j+1}\\s_j&s_{j+1}\end{pmatrix}.
\]
The entry's requirement on that subblock is
\begin{equation}\label{eq:loc}
4-\#\{\beta-\alpha,\ \delta-\beta,\ \gamma-\delta,\ \alpha-\gamma\}\;=\;1,
\end{equation}
a condition on the four entries alone.

For $r,s\in V$ declare $r\to s$ an edge of a digraph $G$ when, for every
$1\le j\le 6-1$,
\begin{itemize}
\item the four entries above satisfy \eqref{eq:loc};

\end{itemize}
Let $M\in\{0,1\}^{S\times S}$ be the adjacency matrix of $G$ and $\mathbf{1}$ the
all-ones vector.

\begin{lemma}\label{lem:walk}
The arrays counted by the entry with $L$ rows are exactly the walks
$r^{(1)}\to r^{(2)}\to\cdots\to r^{(L)}$ of length $L-1$ in $G$. Consequently
\[
a(n)\;=\;\tfrac1{2}\mathbf{1}^{\!\top}M^{n}\mathbf{1} .
\]
\end{lemma}

\begin{proof}
An array with $L$ rows and $6$ columns is precisely a sequence
$r^{(1)},\dots,r^{(L)}$ of elements of $V$; conversely any such sequence is an array.
Every $2\times2$ subblock lies in two consecutive rows $r^{(t)},r^{(t+1)}$ and two
consecutive columns, so the array satisfies the entry's condition on all of its subblocks
if and only if each consecutive pair $\bigl(r^{(t)},r^{(t+1)}\bigr)$ satisfies it for
every column pair --- which is exactly the edge relation of $G$. The number of walks of
length $L-1$, summed over all starting and ending vertices, is
$\mathbf{1}^{\!\top}M^{L-1}\mathbf{1}$, since $(M^{k})_{r,s}$ counts the walks of
length $k$ from $r$ to $s$. The entry's index $n$ corresponds to $L=n+1$ rows. The entry counts $\tfrac1{2}$ of those arrays, a constant factor that passes through every step below unchanged.
\end{proof}

\section{The criterion}

Let $q(t)=t^{70}-\sum_i c_i\,t^{70-i}$ be the characteristic polynomial of
\eqref{eq:conj}, with the $c_i$ read off that equation.

\begin{lemma}\label{lem:crit}
For every $n$ with $n+1-1\ge70$,
\[
a(n)-\sum_i c_i\,a(n-i)\;=\;\tfrac1{2}\mathbf{1}^{\!\top}M^{\,n+1-1-70}\,q(M)\,\mathbf{1} .
\]
Hence \eqref{eq:conj} holds from that point on if and only if
$\mathbf{1}^{\!\top}M^{j}q(M)\mathbf{1}=0$ for every $j\ge0$; and it is enough to check
$j=0,1,\dots,S-1$.
\end{lemma}

\begin{proof}
By Lemma~\ref{lem:walk}, $a(n-i)=\tfrac1{2}\mathbf{1}^{\!\top}M^{\,n+1-1-i}\mathbf{1}$, so
\[
a(n)-\sum_i c_i a(n-i)
=\tfrac1{2}\mathbf{1}^{\!\top}\Bigl(M^{m}-\sum_i c_i M^{m-i}\Bigr)\mathbf{1}
=\tfrac1{2}\mathbf{1}^{\!\top}M^{\,m-70}\,q(M)\,\mathbf{1}
\]
with $m=n+1-1$, which is the stated identity; the scalar $\tfrac1{2}$ never vanishes, so
it does not affect whether the left side is zero. The equivalence follows by letting
$m-70=j$ run over $\ge0$. For the last clause put $w=q(M)\mathbf{1}$. By the
Cayley--Hamilton theorem $M^{S}$ is an integer combination of $I,M,\dots,M^{S-1}$, so
$\mathbf{1}^{\!\top}M^{j}w$ for $j\ge S$ is the corresponding combination of the earlier
values; if those all vanish, so do all the rest.
\end{proof}

\section{The computation}

\begin{theorem}
The recurrence \eqref{eq:conj} holds for every $n>69$, which is the statement of
Section~1.
\end{theorem}

\begin{proof}
The digraph was constructed from \eqref{eq:loc} by a depth-first scan across the $6$
columns: the edge condition is a conjunction of one constraint per consecutive column
pair, so the successors of a row $r$ are enumerated column by column without testing all
$S^2$ pairs. The vector $w=q(M)\mathbf{1}\in\mathbb{Z}^{S}$ was then formed by $70$
matrix--vector products in exact integer arithmetic, and $\mathbf{1}^{\!\top}M^{j}w$ was
evaluated for $j=0,1,\dots$, again exactly. Every one of those integers vanishes from the
index corresponding to $n=69+1$ onwards, and the run continued until $w$ itself was the
zero vector, which settles every larger $j$ at once. By Lemma~\ref{lem:crit} the recurrence
holds for every $n>69$.
\end{proof}

\begin{remark}
The argument gives more than the recurrence: it shows the sequence is $C$-finite with
characteristic polynomial dividing that of $M$, so it has a rational generating function
whose denominator divides $\det(I-xM)$. The conjectured recurrence is one consequence; the
minimal one is a divisor of $q$.
\end{remark}

\section{Verification}

Three checks, each independent of the algebra above.

First, the digraph was built directly from the entry's stated condition and the walk counts
$\tfrac1{2}\mathbf{1}^{\!\top}M^{L-1}\mathbf{1}$ compared against the entry's DATA at
every one of the $11$ published indices; the values agree exactly. That is what ties
the digraph to the sequence: a model that reproduced only some of the published terms would
be the wrong model, and would have been rejected. In particular it is what rules out a
misreading of the entry's wording, since a different reading of the condition gives a
different digraph and different counts.

Second, the recurrence \eqref{eq:conj} was evaluated on those published terms in exact
integer arithmetic, with no matrices involved, and vanishes wherever it applies.

Third, the annihilation test of Lemma~\ref{lem:crit} was run in exact integer arithmetic
until the working vector was identically zero, not to a sampled prefix, and returned zero
throughout.

\begin{thebibliography}{9}
\bibitem{oeis} The OEIS Foundation, \emph{The On-Line Encyclopedia of Integer
Sequences}, \url{https://oeis.org}, 2026. Sequence A209793.
\bibitem{stanley} R.~P.~Stanley, \emph{Enumerative Combinatorics}, Volume 1, 2nd ed.,
Cambridge University Press, 2011. (Transfer-matrix method, Section 4.7.)
\bibitem{fs} P.~Flajolet and R.~Sedgewick, \emph{Analytic Combinatorics}, Cambridge
University Press, 2009.
\bibitem{horn} R.~A.~Horn and C.~R.~Johnson, \emph{Matrix Analysis}, 2nd ed., Cambridge
University Press, 2013.
\end{thebibliography}

\end{document}
